% !TeX program = xelatex
\documentclass[11pt,fleqn]{article}
\usepackage[letterpaper,margin=0.92in]{geometry}
\usepackage{fontspec}
\IfFontExistsTF{TeX Gyre Pagella}
  {\setmainfont{TeX Gyre Pagella}}
  {\setmainfont{TeXGyrePagella}}
\IfFontExistsTF{TeX Gyre Heros}
  {\setsansfont{TeX Gyre Heros}}
  {\setsansfont{TeXGyreHeros}}
\setmonofont{DejaVu Sans Mono}
\usepackage{eulervm}
\usepackage{microtype}
\usepackage{amsmath,amssymb,mathtools,amsthm}
\usepackage{xcolor,hyperref,booktabs,array,enumitem,fancyhdr}
\hypersetup{colorlinks=true,linkcolor=black,urlcolor=blue!55!black,citecolor=black}
\setlength{\parindent}{1.2em}
\setlength{\parskip}{0pt}
\setlength{\mathindent}{1.7em}
\setlist{nosep,leftmargin=1.6em}
\setlength{\headheight}{14pt}
\pagestyle{fancy}
\fancyhf{}
\fancyhead[L]{\small Elliptic Laurent periods and A216584}
\fancyhead[R]{\small Bradley Klee}
\fancyfoot[C]{\thepage}
\renewcommand{\headrulewidth}{0.25pt}
\renewcommand{\footrulewidth}{0pt}
\allowdisplaybreaks
\theoremstyle{plain}
\newtheorem{theorem}{Theorem}
\newtheorem{proposition}{Proposition}
\newcommand{\Dx}{D_x}
\newcommand{\Dy}{D_y}

\title{\vspace{-2.1em}\textbf{Elliptic Laurent periods and A216584}\\[0.3em]
\large with recipes for reduction certificates}
\author{Bradley Klee, \quad Mech.An.ika Sol (Open AI)}
\date{August 8, 2026}

\raggedbottom
\begin{document}
\maketitle

\begin{abstract}
The sequence $a(n)=C_n\;T_n$ where $C_n = [(x+2+x^{-1})^n]_0$ and $T_n = [(y+1+y^{-1})^n]_0$  
gives the product of central binomial coefficients $C_n$ with central trinomial coefficients
$T_n$. The generating function $A(s)=\sum_{n\ge0}a_ns^n$ is also the period function of an
elliptic surface. The ODE constraining $A(s)$ derives from reduction analysis of the 
period integrand and as well from the definition of $C_n$ and $T_n$ as constant terms
of Laurent polynomials. Two approaches yield distinct certificates. 
\end{abstract}

\begin{theorem}[]\label{thm:main}
The sequence $a(n)$ and its generating function $A(s)$ satisfy all of the following:
\begin{enumerate}[label=\textnormal{(\arabic*)},leftmargin=2.2em,itemsep=0.22em,topsep=0.2em]
\item For $G(x,y)=(x+2+x^{-1})(1+y+y^{-1})$,
\[
 a_n=[G(x,y)^n]_0,\qquad
 A(s)=\left[\frac{1}{1-sG(x,y)}\right]_0.
\]

\item Given the elliptic curve $E_s= \{ \; (u,v) : v^2=u^3+(1-4s)u^2+16s^2u \; \}$ then 
\[
 % is k = 1? 
 A(s)=c_\gamma^{-1} T_\gamma(k s) \qquad \text{where} \qquad 
 T(s) = \oint_{\gamma} \frac{du}{v} = \frac{du}{\sqrt{u^3+(1-4s)u^2+16s^2u}}
\]
and $T(s)$ is both an elliptic integral and the period of $E_s$. 

\item The terms $a(n)$ satisfy a p-recurrence: 
\[
 n^2a_n=2(2n-1)^2a_{n-1}
       +12(2n-1)(2n-3)a_{n-2}\qquad(n\ge2),
\]
and the generating function satisfies an ordinary differential equation: 
\[
L \circ A(s) =  s(1-12s)(1+4s)A''+(1-16s-144s^2)A'-(2+36s)A=0.
\]

\item Approaches (1) and (2) arrive at the result (3) through distinct certificates
\[
 L\circ \!\left(\frac{du}{\sqrt P}\right)
   =d_u\!\left(\frac{R}{P^{3/2}}\right),\qquad
 L\circ \!\left(\frac1\rho\right)
   =D_x\!\left(\frac{\Xi_x}{\rho^2}\right)
    +D_y\!\left(\frac{\Xi_y}{\rho^2}\right),
\]
where $\rho = 1-sG$, $D_x=x\partial_x$ and $D_y=y\partial_y$, and numerators are 
easier for machines to read. The exact finite certificate data are recorded in the 
embedded JSON certificates designed for automatic checking.
\end{enumerate}
The initial values are
\[
1,2,18,140,1330,12852,130284,1348776,\ldots
\]
\end{theorem}

\clearpage

\section{A216584 as curve geometry}

\subsection*{Curve certificate and recurrence}
Set
\[
 L=s(1-12s)(1+4s)\partial_s^2
 +(1-16s-144s^2)\partial_s-(2+36s).
\]
Direct differentiation of \(P^{-1/2}\) gives the exact identity
\begin{equation}\label{eq:curve-cert}
 L\!\left(\frac{du}{\sqrt P}\right)
 =d_u\!\left(\frac{R(u,s)}{P^{3/2}}\right),
\end{equation}
where
\begin{equation}\label{eq:Rcurve}
\begin{aligned}
 R(u,s)=4u^2\bigl(&96s^3-48s^2u-16s^2+18su^2+24su+8s\\
                 &{}+u^2+u\bigr).
\end{aligned}
\end{equation}
After clearing \(P^{5/2}\), both sides of \eqref{eq:curve-cert} have
identical polynomial numerator.  Hence every closed period of \(E_s\) satisfies \(L\Pi=0\).

Substitution of \(A(s)=\sum a_ns^n\) into the differential equation in
Theorem~\ref{thm:main}(3) gives the recurrence.  With \(a_0=1\) and \(a_1=2\), it begins
\[
 1,2,18,140,1330,12852,130284,1348776,\ldots .
\]

\subsection*{Product factorization and Bala's diagonal}
The central trinomial coefficients satisfy
\[
 nT_n=(2n-1)T_{n-1}+3(n-1)T_{n-2},
\]
while
\[
 \frac{\binom{2n}{n}}{\binom{2n-2}{n-1}}
 =\frac{2(2n-1)}{n}.
\]
Their product therefore satisfies Theorem~\ref{thm:main}(3); the initial values
identify
\begin{equation}\label{eq:hadamard}
 a_n=\binom{2n}{n}T_n.
\end{equation}
The factorization $G$ makes this identity literal at the Laurent
level:
\[
 [G^n]_0
 =[(x+2+x^{-1})^n]_0\,
  [(1+y+y^{-1})^n]_0
 =\binom{2n}{n}T_n.
\]

Peter Bala recorded on OEIS the diagonal representation
\begin{equation}\label{eq:bala}
 a_n=[X^nY^n]\bigl(1+(X+Y)^2+(X+Y)^4\bigr)^n.
\end{equation}
To see the same factorization directly, write
\[
 (1+z+z^2)^n=\sum_k c_{n,k}z^k.
\]
In \eqref{eq:bala}, only \(k=n\) can contribute to total bidegree
\((n,n)\).  Hence
\[
 [X^nY^n](\cdots)^n
 =c_{n,n}[X^nY^n](X+Y)^{2n}
 =T_n\binom{2n}{n}.
\]
Thus $G$ is a factorized constant-term form of the same known
product structure.

\subsection*{The Laurent pencil and the elliptic curve}
The level equation \(1-sG=0\) is
\begin{equation}\label{eq:G-pencil}
 s(x+1)^2(y^2+y+1)-xy=0.
\end{equation}
Define
\begin{equation}\label{eq:G-to-E}
 u=-4sy,
 \qquad
 v=4sy\,\frac{x-1}{x+1}.
\end{equation}
A direct substitution gives
\[
 v^2-u^3-(1-4s)u^2-16s^2u
 =\frac{64s^2y}{(x+1)^2}
 \bigl(s(x+1)^2(y^2+y+1)-xy\bigr).
\]
Thus \eqref{eq:G-to-E} maps the Laurent pencil birationally to $E_s$;
away from the exceptional divisors the inverse is
\begin{equation}\label{eq:E-to-G}
 y=-\frac{u}{4s},
 \qquad
 x=\frac{u-v}{u+v}.
\end{equation}
The period differential is preserved as well.  If
\(Q=s(x+1)^2(y^2+y+1)-xy\), then
\[
 \frac{dx/x\wedge dy/y}{1-sG}=-\frac{dx\wedge dy}{Q}.
\]
On \(Q=0\),
\[
 Q_x=y\frac{x-1}{x+1},
\]
and therefore
\begin{equation}\label{eq:residue-map}
 \operatorname{Res}_{Q=0}\frac{dx/x\wedge dy/y}{1-sG}
 =-\frac{dy}{Q_x}
 =\frac{du}{v}.
\end{equation}
This proves the period identity in Theorem~\ref{thm:main}(2) directly, not merely
by comparison of differential equations.

For orientation, the discriminant of $E_s$ is
\[
 -4096s^4(1+4s)(12s-1).
\]
The minimal rational elliptic surface has fibres
\(I_4,I_1,I_1,I_0^*\) at \(0,-1/4,1/12,\infty\), respectively.  The visible
2-torsion point \((0,0)\) gives the 2-isogenous companion
\[
 v^2=u^3+(8s-2)u^2+(1-12s)(1+4s)u,
\]
whose exact differential-equation certificate is included in the accompanying
verification file; it has the same differential operator as Theorem~\ref{thm:main}(3).

The same second-order equation also controls the other local periods of this
surface.  In particular, with
\[
 H(z)={}_2F_1\!\left(\frac12,\frac12;1;z\right),
\]
one may take local solutions
\[
 \Pi_{1/12}(q)=\sqrt{1+q/3}\,H(q),
 \qquad q=\frac{1-12s}{1+4s},
\]
and
\[
 \Pi_{-1/4}(r)=\sqrt{1-4r}\,H(r),
 \qquad r=\frac{1+4s}{16s}.
\]
The A216584 branch is the analytic solution normalized at \(s=0\); the
second local solution there is logarithmic.

\section{The Laurent--ODE first approach}

We now forget the curve, the termwise product factorization, and the known recurrence,
and start only from \(G\).  Write
\[
 \rho=1-sG(x,y),
 \qquad
 \Phi(s)=\left[\frac1\rho\right]_0.
\]
Equivalently,
\begin{equation}\label{eq:torus-period}
 \Phi(s)=\frac{1}{(2\pi i)^2}
 \oint\!\oint
 \frac{dx}{x}\frac{dy}{y}\frac1\rho .
\end{equation}
Let \(\theta=s\partial_s\).  Successive differentiated integrands are
\[
 \frac1\rho,
 \qquad
 \theta\frac1\rho=\frac{sG}{\rho^2},
 \qquad
 \theta^2\frac1\rho
 =\frac{sG}{\rho^2}+\frac{2s^2G^2}{\rho^3}.
\]
The reduction asks whether these classes become linearly dependent modulo
\(\Dx\)- and \(\Dy\)-divergences.

For this input one may take
\[
 B=\{x^iy^j:-2\le i,j\le2\},
 \qquad
 C=\{x^ky^\ell:-3\le k,\ell\le3\}.
\]
Thus \(|B|=25\) and \(|C|=49\).  Here the basis-enlargement parameter is
\(d=2\): explicitly,
\[
 B=(2\,\operatorname{conv}(\{0\}\cup\operatorname{supp}G))\cap\mathbb Z^2
   =[-2,2]^2\cap\mathbb Z^2.
\]
Since \(\operatorname{supp}G=[-1,1]^2\cap\mathbb Z^2\), multiplication by
\(\rho=1-sG\) enlarges the exponent region by one unit in each direction,
so
\[
 C=[-3,3]^2\cap\mathbb Z^2.
\]
This \(d\) is a support bound and is independent of the pole-order parameter
below.  No support outside \(C\) can occur in the linear system for this
chosen ansatz.

The matrices are large only in dimension; their entries are given by a small
fixed pattern of exponent shifts.  Write
\[
 f_{r,q}=[x^ry^q]G,
\]
with \(f_{r,q}=0\) outside \(\{-1,0,1\}^2\).  For
\(\alpha=(k,\ell)\in C\) and \(\beta=(i,j)\in B\), the three \(25\)-column
blocks of \(G\) are
\begin{equation}\label{eq:G-entries}
\begin{aligned}
 G^{(0)}_{\alpha,\beta}
   &=\delta_{\alpha,\beta}-s f_{k-i,\ell-j},\\
 G^{(x)}_{\alpha,\beta}
   &=s(k-i)f_{k-i,\ell-j},\\
 G^{(y)}_{\alpha,\beta}
   &=s(\ell-j)f_{k-i,\ell-j}.
\end{aligned}
\end{equation}
Hence every column of the first block has at most nine nonzero entries and
every column of either derivative block has at most six.  In all,
\[
 G:49\times75,\qquad \#\operatorname{nz}(G)=525,
\]
so only \(14.3\%\) of its entries are nonzero.  The divergence matrix is even
simpler: for \(\beta=(i,j)\),
\begin{equation}\label{eq:J-entries}
 J^{(x)}_{\beta,\gamma}=i\,\delta_{\beta,\gamma},
 \qquad
 J^{(y)}_{\beta,\gamma}=j\,\delta_{\beta,\gamma}.
\end{equation}
Thus \(J:25\times50\) has only 40 nonzero entries.  Because a second
\(s\)-derivative produces a numerator over \(\rho^3\), only one pole-lowering
step is needed, with pole-order parameter \(r=2\).  The complete fixed
order-two system is
\[
 \begin{pmatrix}
   G&0&0\\
   I_{25}&-J/2&-[\rho]_B&-[G]_B
 \end{pmatrix}
 \begin{pmatrix}U\\V_x\\V_y\\\alpha\\\beta\end{pmatrix}
 =
 \begin{pmatrix}E[G^2]_B\\0\end{pmatrix}.
\]
It has size \(74\times77\) and only 608 nonzero entries (\(10.7\%\)).

For this example the rational reconstruction is also finitely bounded.  The
chosen exact solution has numerator degree at most 3 and denominator degree at
most 4, hence total degree at most 7.  Total-degree caps 5 and 6 fail, whereas
7 succeeds and passes complete symbolic replay.  Accordingly the fixed
reconstruction used below takes eight regular integer values to determine the
rational functions and three further values, not used in that determination,
as independent checks.  After interpolation every
entry is substituted back into the original \(74\times77\) system over
\(\mathbb Q(s)\); the certificate does not rest on the sampled equations.

The resulting differential relation is
\begin{equation}\label{eq:theta-op}
  A_\theta=
 \theta^2-2s(2\theta+1)^2
 -12s^2(2\theta+1)(2\theta+3).
\end{equation}
In ordinary derivatives, \eqref{eq:theta-op} is exactly
Theorem~\ref{thm:main}(3).

The same solve returns the certificate rather than requiring a second
calculation.  There are Laurent polynomials \(\Xi_x,\Xi_y\), supported on the
25 exponent pairs of \(B\), such that
\begin{equation}\label{eq:Xi-identity}
 A\!\left(\frac1\rho\right)
 =\Dx\!\left(\frac{\Xi_x}{\rho^2}\right)
 +\Dy\!\left(\frac{\Xi_y}{\rho^2}\right),
\end{equation}
where
\[
 A=(-48s^3-8s^2+s)\partial_s^2
   +(-144s^2-16s+1)\partial_s-(36s+2).
\]
This is the same operator as Theorem~\ref{thm:main}(3).  The exact sparse
coefficients of \(\Xi_x,\Xi_y\), together with the ordered bases and matrices,
are stored in
\path{examples/public/low_support/certificates/A216584_G_guvj.json}.
The publication replay checks, by coefficient comparison,
\[
 \begin{aligned}
 GUVJ&=0, & Xp&=0,\\
 A(1/\rho)-\Dx(\Xi_x/\rho^2)-\Dy(\Xi_y/\rho^2)&=0.
 \end{aligned}
\]
It also regenerates the coefficients and verifies Theorem~\ref{thm:main}(3).
All five exact checks pass.  No finite-term fit is used to obtain
\eqref{eq:theta-op}.

Thus the curve calculation of Section~1 and the Laurent reduction of this
section meet at the same certified differential equation by independent
routes.

\section{Formal reductions}

The preceding calculation generalizes.  Let
\(F\in K[x^{\pm1},y^{\pm1}]\), write
\[
 \rho=1-tF,
 \qquad
 \Dx=x\frac{\partial}{\partial x},
 \qquad
 \Dy=y\frac{\partial}{\partial y},
\]
and consider
\[
 \Phi_F(t)=\left[\frac1\rho\right]_0.
\]
The constant term of either \(\Dx(R)\) or \(\Dy(R)\) is zero.  Hence an
identity
\begin{equation}\label{eq:general-cert}
 A\!\left(\frac1\rho\right)=\Dx(\Xi_x)+\Dy(\Xi_y)
\end{equation}
proves \(A\Phi_F=0\).

Fix a pole order \(r\).  If Laurent polynomials \(w,a,b_x,b_y\) satisfy
\begin{equation}\label{eq:split}
 w=\rho a-\Dx(\rho)b_x-\Dy(\rho)b_y,
\end{equation}
then
\begin{equation}\label{eq:lower}
 \frac{w}{\rho^{r+1}}
 =\frac{a-r^{-1}(\Dx b_x+\Dy b_y)}{\rho^r}
 +\Dx\!\left(\frac{b_x}{r\rho^r}\right)
 +\Dy\!\left(\frac{b_y}{r\rho^r}\right).
\end{equation}
This is the elementary pole-lowering identity behind the calculation.

Let \(P=\operatorname{conv}(\{0\}\cup\operatorname{supp}F)\).  Choose an
integer basis-enlargement parameter \(d\ge1\) and take
\[
 B_d=(dP)\cap\mathbb Z^2.
\]
Let \(C_d\) be the union of \(B_d\) with all translates
\(B_d+u\), \(u\in\operatorname{supp}F\).  These are explicit finite exponent
bounds: if \(a,b_x,b_y\) are supported in \(B_d\), every coefficient equation
in \eqref{eq:split} lies in \(C_d\).  The parameter \(d\) controls only this
support ansatz; it is distinct from the pole order \(r\) in
\eqref{eq:lower}.  A successful exact replay proves that the selected bound is
sufficient for that certificate.  We make no claim here that a fixed \(d\)
works for every Laurent polynomial.

Let \(B=(b_1,\ldots,b_N)\) be the ordered monomial basis associated with
\(B_d\), and let \(C\) be the ordered monomial basis associated with \(C_d\).
The coefficient matrix
\[
 G\in K(t)^{M\times3N}
\]
represents the map
\[
 (a,b_x,b_y)\longmapsto
 \rho a-\Dx(\rho)b_x-\Dy(\rho)b_y.
\]
If \(E\) includes the \(B\)-space into the \(C\)-space, exact solution of
\begin{equation}\label{eq:GUV}
 G\begin{pmatrix}U\\V_x\\V_y\end{pmatrix}=E
\end{equation}
produces both the reduced part and the divergence data.  The matrix
\[
 J: (b_x,b_y)\longmapsto \Dx(b_x)+\Dy(b_y)
\]
then gives the one-step reduction
\begin{equation}\label{eq:lower-matrix}
 \operatorname{Lower}_r=U-\frac1rJ
 \begin{pmatrix}V_x\\V_y\end{pmatrix}.
\end{equation}
Linear dependence among finitely many reduced derivative columns gives the
differential operator \(A\); the same linear combination of the stored \(V\)-data gives
the exact certificate \(\Xi\).  The basis-enlargement parameter is independent
of the pole order \(r\); the implementation records both explicitly.

For comparison, the bounded low-support search that led back to A216584 also
produced the following small certified OEIS hits:
\begin{center}
\small
\begin{tabular}{p{0.15\textwidth}p{0.48\textwidth}p{0.12\textwidth}p{0.12\textwidth}}
\toprule
OEIS & Laurent polynomial & order & exact certificate\\
\midrule
A002894 & \(x^2+xy+(xy)^{-1}+x^{-2}\) & 2 & yes\\
A006480 & \(xy+x/y+x^{-2}\) & 2 & yes\\
A216584 & \((x+2+x^{-1})(1+y+y^{-1})\) & 2 & yes\\
\bottomrule
\end{tabular}
\end{center}
The first two are retained as regression cases rather than developed here.
The same bounded search also produces additional period fingerprints with no
OEIS identification in the initial lookup; those are candidates for later
geometric study, not claims of new sequences in this paper.

\section{A bounded replay for A216584}\label{sec:pseudocode}

The general theory above admits many implementation choices.  For the main
example we remove those choices.  The following pseudocode is a specification
of one finite calculation, not a heuristic search procedure.  It fixes the
support enlargement, differential order, pole order, interpolation bound, and
validation count in advance.  An independent implementation therefore has no
operator-order search or basis search to reproduce.

\begingroup
\fontsize{9.25}{10.65}\selectfont
\begin{enumerate}[label=\texttt{\arabic*.},leftmargin=2.7em,itemsep=0.25em,topsep=0pt]
\item Set
\[
 F=(x+2+x^{-1})(1+y+y^{-1}),\qquad \rho=1-sF,
\]
and use exact rational arithmetic throughout.

\item Fix the support enlargement \(d=2\).  Order lexicographically
\[
 B=[-2,2]^2\cap\mathbb Z^2,\qquad
 C=[-3,3]^2\cap\mathbb Z^2.
\]
Thus \(|B|=25\), \(|C|=49\).  Store a Laurent polynomial as its coefficient
vector on one of these ordered sets.

\item Build \(G\in\mathbb Q(s)^{49\times75}\) from
\[
 (a,b_x,b_y)\longmapsto \rho a-\Dx(\rho)b_x-\Dy(\rho)b_y
\]
using the entry formulas \eqref{eq:G-entries}.  Build the inclusion matrix
\(E\in\mathbb Q^{49\times25}\) and
\(J\in\mathbb Q^{25\times50}\) from \eqref{eq:J-entries}.  Check
\[
 \#\operatorname{nz}(G)=525,\qquad
 \#\operatorname{nz}(J)=40.
\]

\item Form the \(B\)-coordinate columns \(f=[F]_B\),
\(f_2=[F^2]_B\), and \(r_0=[\rho]_B\).  Fix differential order two and pole
order two.  Assemble
\[
 M=
 \begin{pmatrix}
   G&0&0\\
   I_{25}&-J/2&-r_0&-f
 \end{pmatrix},\qquad
 b=\begin{pmatrix}Ef_2\\0\end{pmatrix}.
\]
Check \(M\) has shape \(74\times77\) and 608 nonzero entries.  The unknown is
\(z=(U,V_x,V_y,\alpha,\beta)^T\).

\item At the first regular positive integer value of \(s\), determine pivot
columns and an equal number of independent rows.  Keep that square pivot block
fixed at every later sample.  Set all free variables to zero.

\item Fix the rational reconstruction bound
\[
 \deg(\text{numerator})+\deg(\text{denominator})\le7.
\]
Solve the fixed square block exactly at regular integer values of \(s\).  Use
eight values to determine the rational functions and three additional regular
values only as checks.  For each pivot variable reconstruct the least-total-
degree rational function satisfying the first eight values and reject it unless
it also satisfies all three check values.

\item Reassemble the symbolic vector \(z\in\mathbb Q(s)^{77}\) and require the
full exact identity
\[
 Mz=b.
\]
This symbolic replay, not interpolation, is the proof of the reconstructed
solution.  Split \(z\) into \(U,V_x,V_y,\alpha,\beta\).

\item Verify separately
\[
 G\begin{pmatrix}U\\V_x\\V_y\end{pmatrix}=Ef_2,
 \qquad
 U-\frac12J\begin{pmatrix}V_x\\V_y\end{pmatrix}
   =\alpha r_0+\beta f.
\]
These are the \(G,U,V,J\) reduction and the single pole-lowering closure.

\item Form
\[
 X=\begin{pmatrix}1&0&2\alpha\\0&1&2\beta\end{pmatrix}.
\]
Take the primitive polynomial kernel vector
\(p=(p_0,p_1,p_2)^T\), normalized by denominator clearing, polynomial-gcd
removal, integer-content removal, and sign fixing.  Set
\[
 A=p_0+p_1\partial_s+p_2\partial_s^2.
\]

\item Set \(V=(V_x,V_y)^T\) and \(\Xi=p_2V\).  Verify the cleared coefficient
identity equivalent to
\[
 A(1/\rho)
 =\Dx(\Xi_x/\rho^2)+\Dy(\Xi_y/\rho^2).
\]
Require every coefficient in the ordered Laurent basis to vanish exactly.

\item Convert \(A\) to \(\theta=s\partial_s\) form and obtain the recurrence by
replacing each \(s^kQ_k(\theta)\) with
\(Q_k(n-k)a_{n-k}\).  The result must be Theorem~\ref{thm:main}(3).

\item Generate twelve constant terms by exact lattice convolution, replay the
recurrence on every available index, and independently compute the first five
constant terms by direct expansion.  Require both lists to agree.

\item Serialize \(B,C,G,J,U,V_x,V_y,\alpha,\beta,A,\Xi\), the recurrence,
and the term checks.  A verifier must rebuild the coefficient identities from
these data rather than trust stored pass flags.
\end{enumerate}
\endgroup

For the exact solution produced by this replay, the largest numerator degree is
3 and the largest denominator degree is 4.  The degree-7 reconstruction bound
is therefore attained; repeating the same reconstruction with caps 5 or 6
fails, while cap 7 succeeds and the resulting functions pass the full symbolic
system.  These numerical degree bounds concern reconstruction in \(s\); they
are separate from both the exponent bound \(d=2\) and the pole order \(r=2\).

The accompanying implementation exposes the same calculation through
\begin{verbatim}
python3 code/guvj/guvj_period_factory.py \
  --F "(x + 2 + 1/x)*(1 + y + 1/y)" \
  --model A216584_G \
  --dilation 2 \
  --max-total-degree 7 \
  --output examples/public/low_support/certificates/A216584_G_guvj.json

python3 code/public/verify_a216584_geometry.py
\end{verbatim}
The first command builds the Laurent certificate.  The second independently
checks the elliptic differential identity, recurrence, termwise product, Bala
diagonal formula, birational map, residue differential, and agreement with the
stored constant terms.  The paper deliberately specifies only the A216584
instance in algorithmic detail; Section~3 gives the general reduction from
which it is obtained.

\section*{Computational disclosure}
The symbolic calculations and certificate replays in this paper were performed
with the accompanying Python/SymPy implementation.  OpenAI language models
were used during development for mathematical exploration, code assistance,
testing, and editorial work.  Claims in the paper do not rely on uncheckable
model output: the asserted differential identities, recurrences, birational
maps, and Laurent reductions are represented by exact symbolic identities or
machine-readable certificates that can be replayed independently.

\begin{thebibliography}{9}

\bibitem{A216584}
OEIS Foundation Inc., \emph{The On-Line Encyclopedia of Integer Sequences},
A216584.  The entry includes the product formula of Paul D. Hanna, Peter
Bala's diagonal formula (2020), and Mark van Hoeij's hypergeometric generating
function (2025).

\bibitem{A002894}
OEIS Foundation Inc., \emph{The On-Line Encyclopedia of Integer Sequences},
A002894.

\bibitem{A006480}
OEIS Foundation Inc., \emph{The On-Line Encyclopedia of Integer Sequences},
A006480.

\bibitem{Klee}
B. J. Klee,
\emph{An Update on the Computational Theory of Hamiltonian Period Functions},
Ph.D. dissertation, University of Arkansas, Fayetteville, 2020, especially
Chapter~3.

\bibitem{BDS}
A. Bostan, L. Dumont, and B. Salvy,
``Algebraic Diagonals and Walks: Algorithms, Bounds, Complexity,''
\emph{Journal of Symbolic Computation} 83 (2017), 68--92.

\bibitem{BCChL}
A. Bostan, S. Chen, F. Chyzak, and Z. Li,
``Complexity of creative telescoping for bivariate rational functions,''
\emph{Proceedings of ISSAC 2010}, 203--210.

\bibitem{AZ}
G. Almkvist and D. Zeilberger,
``The method of differentiating under the integral sign,''
\emph{Journal of Symbolic Computation} 10 (1990), 571--591.

\bibitem{SymPy}
A. Meurer et al.,
``SymPy: symbolic computing in Python,''
\emph{PeerJ Computer Science} 3 (2017), e103.

\end{thebibliography}


\end{document}
